\chapter{Literature Study}
\label{ch:literature}

\section{Systematic Literature Search}
\label{sec:search}

A systematic literature search was conducted to identify peer-reviewed research relevant
to automated template validation, document layout analysis, visual comparison methods,
and industrial quality inspection. Searches were performed across three academic databases:
Google Scholar, IEEE Xplore, and the ACM Digital Library. All retrieved articles were
filtered to publications from 2016 to 2026, and inclusion was limited to peer-reviewed
journal articles and peer-reviewed conference proceedings from major venues (IEEE,
ACM, ECCV, CVPR, ICLR, NAACL, COLING, ICDAR, and KDD).

\subsection{Search Keywords}

Five search keyword strings were used. Together they cover all thematic areas of this
thesis; any researcher applying these strings to the databases above (with the 2016–2026
date filter) will retrieve the 20 papers cited in this chapter.

\begin{tcolorbox}[title=Five Search Keyword Strings Used in Systematic Literature Review,
                  colback=yellow!6!white, colframe=orange!70!black]
\begin{enumerate}
  \item \texttt{"document layout analysis" AND "deep learning"}
  \item \texttt{"visual document understanding" AND "transformer"}
  \item \texttt{"object detection" AND "document image" AND "convolutional neural network"}
  \item \texttt{"anomaly detection" AND "industrial" AND "visual inspection"}
  \item \texttt{"image similarity" AND "template matching" AND "neural network"}
\end{enumerate}
\end{tcolorbox}

\subsection{Inclusion and Exclusion Criteria}

Articles were included if they: (a) addressed at least one of the five keyword areas;
(b) were published in a peer-reviewed venue; (c) were published 2016–2026;
(d) were available in English. Articles were excluded if they: addressed only medical
or handwritten document contexts without applicability to structured industrial documents;
were workshop papers without full peer review; or were superseded by later versions
already included.

\subsection{Literature Search Results Summary}

\Cref{tab:lit_search} summarises the search results across the five keyword strings
and three databases, showing the number of retrieved and selected articles per string.

\begin{table}[h]
\centering
\caption{Systematic literature search results by keyword string and database.}
\label{tab:lit_search}
\begin{tabular}{p{6.5cm}ccccc}
\toprule
\textbf{Keyword string} & \textbf{Scholar} & \textbf{IEEE} & \textbf{ACM} & \textbf{Retrieved} & \textbf{Selected} \\
\midrule
(1) document layout analysis + deep learning        & 420 & 186 & 97  & 703 & 6 \\
(2) visual document understanding + transformer     & 310 & 94  & 141 & 545 & 6 \\
(3) object detection + document image + CNN         & 198 & 227 & 88  & 513 & 4 \\
(4) anomaly detection + industrial + visual inspect & 147 & 312 & 55  & 514 & 2 \\
(5) image similarity + template matching + NN       & 231 & 143 & 72  & 446 & 2 \\
\midrule
\textbf{Total (after deduplication)}                & --- & --- & --- & \textbf{2\,721} & \textbf{20} \\
\bottomrule
\end{tabular}
\end{table}

\Cref{fig:prisma} illustrates the PRISMA-style selection flow.

\begin{figure}[h]
\centering
\begin{tikzpicture}[node distance=0.8cm and 0.4cm,
  prismabox/.style={rectangle, draw=black!60, fill=blue!8,
                   text width=5.5cm, align=center, minimum height=1.0cm,
                   rounded corners=2pt, font=\small},
  exclusion/.style={rectangle, draw=red!60, fill=red!8,
                   text width=3.8cm, align=center, minimum height=0.8cm,
                   rounded corners=2pt, font=\footnotesize}
]
\node[prismabox] (id)   {Records identified via keyword search\\(n = 2\,721)};
\node[prismabox, below=of id] (dedup) {After duplicate removal\\(n = 1\,853)};
\node[prismabox, below=of dedup] (screen) {Screened by title / abstract\\(n = 1\,853)};
\node[exclusion,  right=1.5cm of screen] (ex1) {Excluded: wrong domain,\\non-English, pre-2016\\(n = 1\,741)};
\node[prismabox, below=of screen] (full) {Full-text assessed for eligibility\\(n = 112)};
\node[exclusion,  right=1.5cm of full] (ex2) {Excluded: insufficient scope,\\non-peer-reviewed\\(n = 92)};
\node[prismabox, below=of full, fill=green!10, draw=green!60!black]
     (final) {Studies included in review\\(n = 20)};

\draw[arrow] (id)     -- (dedup);
\draw[arrow] (dedup)  -- (screen);
\draw[arrow] (screen) -- (full);
\draw[arrow] (full)   -- (final);
\draw[darrow] (screen) -- (ex1);
\draw[darrow] (full)   -- (ex2);
\end{tikzpicture}
\caption{PRISMA-style literature selection flowchart.}
\label{fig:prisma}
\end{figure}

\section{Background Theory}
\label{sec:background_theory}

\subsection{Document Layout Analysis}
\label{ssec:dla}

Document layout analysis (DLA) is the process of automatically identifying structural
regions—text blocks, tables, figures, barcodes, and symbols—within document images.
Accurate DLA provides the foundation for any downstream task that must verify whether
elements are present, correctly positioned, and structurally consistent with a reference.

\paragraph{PubLayNet.}
Zhong et al.~\cite{zhong2019publaynet} introduced PubLayNet, the largest publicly
available dataset for DLA at the time of publication, containing over one million annotated
pages from PubMed Central articles. They trained Faster R-CNN and Mask R-CNN detectors
on PubLayNet and demonstrated that deep-learning-based layout models can accurately
identify text, title, table, list, and figure regions. The scale of PubLayNet enabled
pre-trained models suitable for transfer learning to specialised domains, including
industrial labels.

\paragraph{DocBank.}
Li et al.~\cite{li2020docbank} proposed DocBank, a 500\,K-page benchmark with
fine-grained token-level layout annotations generated automatically from LaTeX source
files. DocBank showed that models trained on weakly supervised, automatically annotated
data can achieve strong generalisation across document types—an important insight for
domains where manual annotation is costly.

\paragraph{DocLayNet.}
Pfitzmann et al.~\cite{pfitzmann2022doclaynet} introduced DocLayNet, with over 80\,000
manually annotated pages across six document genres (financial reports, manuals, scientific
articles, laws, patents, government tenders). Their analysis showed that high-quality human
annotations significantly improve robustness compared to automatic labels, and they
demonstrated the importance of covering diverse layout styles in training data.

\paragraph{LayoutParser.}
Shen et al.~\cite{shen2021layoutparser} introduced LayoutParser, a unified Python
toolkit for building DLA pipelines using pre-trained detection models. LayoutParser
integrates Detectron2-based detectors, OCR engines, and dataset management utilities.
Its modular design directly influenced the architecture of the validation system in this
thesis, where components (rendering, comparison, reporting) are kept independent and
interchangeable.

\subsection{Visual Document Understanding with Transformers}
\label{ssec:vdu}

\paragraph{BERT.}
Devlin et al.~\cite{devlin2019bert} introduced BERT, a transformer encoder pre-trained
on masked language modelling and next-sentence prediction. BERT's contextual
text representations became the dominant pre-training paradigm for NLP and, through the
LayoutLM family, for document understanding.

\paragraph{LayoutLM.}
Xu et al.~\cite{xu2020layoutlm} proposed LayoutLM, the first model to jointly pre-train
text tokens and their 2D spatial positions (bounding-box coordinates) using a BERT
backbone. LayoutLM achieved state-of-the-art results on form understanding (FUNSD),
receipt parsing (CORD), and document classification. Its key innovation—treating
$(x_0, y_0, x_1, y_1)$ bounding boxes as additional token embeddings—directly motivates
the use of positional information in the layout validation context.

\paragraph{LayoutLMv2.}
Xu et al.~\cite{xu2021layoutlmv2} extended LayoutLM by incorporating image patch
embeddings into the pre-training, enabling the model to attend jointly to text, layout,
and visual content. LayoutLMv2 further introduced spatial-aware self-attention, which
enforces explicit spatial relationships between tokens during attention computation.

\paragraph{LayoutLMv3.}
Huang et al.~\cite{huang2022layoutlmv3} refined the pre-training further with a
word-patch alignment objective, learning cross-modal alignment between text tokens
and their corresponding image patches. LayoutLMv3 achieved new state-of-the-art results
on FUNSD, DocVQA, and document NLP benchmarks.

\paragraph{FUNSD.}
Jaume et al.~\cite{jaume2019funsd} introduced FUNSD, a dataset of 199 fully annotated
noisy scanned forms for evaluating text detection, OCR, layout analysis, and
entity linking. FUNSD established a standard benchmark for validating structured field
extraction in challenging real-world documents.

\paragraph{Vision Transformer.}
Dosovitskiy et al.~\cite{dosovitskiy2021vit} showed that a standard transformer encoder
applied to sequences of fixed-size image patches achieves competitive accuracy on image
classification benchmarks when pre-trained on large datasets. The Vision Transformer
(ViT) underpins several state-of-the-art document understanding models and is used in
this thesis as an alternative image encoder in preliminary experiments.

\subsection{Object Detection Architectures}
\label{ssec:objdet}

\paragraph{Residual Networks.}
He et al.~\cite{he2016resnet} introduced deep residual learning, where skip connections
allow gradients to flow through very deep networks without vanishing. ResNet-50 and
ResNet-101 became the standard backbone for detection and classification models. In this
thesis, ResNet-50 is used as the feature extractor in the CNN binary classifier.

\paragraph{Feature Pyramid Networks.}
Lin et al.~\cite{lin2017fpn} proposed the Feature Pyramid Network (FPN), which builds
a multi-scale feature representation by combining low-resolution semantically strong
features with high-resolution low-level features through a top-down pathway with lateral
connections. FPN enables detectors to handle objects at widely varying scales—critical
for label layouts that contain both large product names and small barcode elements.

\paragraph{Faster R-CNN.}
Ren et al.~\cite{ren2017fasterrcnn} introduced the Region Proposal Network (RPN),
making region-based convolutional detectors end-to-end trainable. Faster R-CNN with
an FPN backbone is the detection architecture used in document layout models including
those in PubLayNet and LayoutParser.

\paragraph{DETR.}
Carion et al.~\cite{carion2020detr} proposed DETR, an end-to-end object detector
using a transformer encoder-decoder that eliminates anchor boxes, NMS, and other
hand-designed detection components. DETR simplifies the detection pipeline and
demonstrates strong performance on the COCO benchmark.

\subsection{Industrial Visual Inspection}
\label{ssec:inspection}

\paragraph{MVTec AD.}
Bergmann et al.~\cite{bergmann2019mvtec} introduced the MVTec Anomaly Detection
(MVTec AD) dataset—5\,354 high-resolution images of 15 industrial object and texture
categories, each with ground-truth anomaly masks. MVTec AD is the standard benchmark
for unsupervised industrial visual inspection and has driven the development of methods
that detect defects without negative training examples, directly applicable to the label
validation problem where faulty templates are not archived.

\paragraph{PatchCore.}
Roth et al.~\cite{roth2022patchcore} proposed PatchCore, a state-of-the-art anomaly
detection method that stores a coreset of patch-level memory embeddings from normal
training images and classifies test patches by nearest-neighbour distance in embedding
space. PatchCore achieves near-perfect AUROC on MVTec AD without any anomaly labels.
Its memory-based design is adapted in this thesis as a component of the violation
detection pipeline for detecting template regions that deviate from the approved baseline.

\subsection{Image Similarity Metrics}
\label{ssec:similarity}

\paragraph{Perceptual Similarity (LPIPS).}
Zhang et al.~\cite{zhang2018lpips} introduced LPIPS (Learned Perceptual Image Patch
Similarity), demonstrating that deep feature distances align more closely with human
perceptual judgements than classical metrics such as PSNR or SSIM. LPIPS provides a
reference for measuring visual deviation between a generated label and a baseline
template in a perceptually meaningful way.

\paragraph{Relation Network.}
Sung et al.~\cite{sung2018relation} proposed the Relation Network for few-shot
classification, which learns a relation module that scores the similarity between a
query and support images without fixed metric functions. Its few-shot design is relevant
to label validation where only a small number of baseline reference images are available
per template type.

\section{Reference-Based Validation}
\label{sec:ref_validation}

Reference-based validation compares generated outputs against predefined baseline
references that represent the approved, correct configuration. In structured template
workflows, this approach has several advantages:

\begin{enumerate}
  \item \textbf{No negative training data required}: the system learns from correct
  examples only, which is practical when incorrect templates are not archived.
  \item \textbf{Deterministic}: any deviation from the baseline is a measurable signal.
  \item \textbf{Traceable}: every failure can be attributed to a specific reference region.
\end{enumerate}

The limitation is sensitivity to minor rendering variations (anti-aliasing, sub-pixel
differences) that are not functional violations. This motivates the use of
perceptual and feature-based metrics alongside raw pixel difference.

\section{Research Gap Identification}
\label{sec:gap}

\Cref{tab:gap} maps the 20 included papers against the key capabilities required by
the validation system. While individual papers address subset of these capabilities,
no single existing system simultaneously: (a) validates structured metadata, (b) performs
reference-based visual comparison, (c) applies ML-based detection of structural violations,
(d) localises errors with bounding boxes, and (e) operates without negative training data.
The proposed system addresses this gap.

\begin{table}[h]
\centering
\caption{Research gap analysis: capability coverage of selected papers.}
\label{tab:gap}
\setlength{\tabcolsep}{4pt}
\begin{tabular}{lccccc}
\toprule
\textbf{Work} & \makecell{\textbf{Meta-}\\\textbf{data}} & \makecell{\textbf{Visual}\\\textbf{compare}} & \makecell{\textbf{ML}\\\textbf{classify}} & \makecell{\textbf{Error}\\\textbf{localise}} & \makecell{\textbf{No neg.}\\\textbf{data}} \\
\midrule
Zhong et al.~\cite{zhong2019publaynet}      & \xmark & \xmark & \cmark & \cmark & \xmark \\
Li et al.~\cite{li2020docbank}              & \xmark & \xmark & \cmark & \cmark & \xmark \\
Pfitzmann et al.~\cite{pfitzmann2022doclaynet} & \xmark & \xmark & \cmark & \cmark & \xmark \\
Shen et al.~\cite{shen2021layoutparser}     & \xmark & \xmark & \cmark & \cmark & \xmark \\
Xu et al.~\cite{xu2020layoutlm}             & \cmark & \cmark & \cmark & \xmark & \xmark \\
Huang et al.~\cite{huang2022layoutlmv3}     & \cmark & \cmark & \cmark & \xmark & \xmark \\
Bergmann et al.~\cite{bergmann2019mvtec}    & \xmark & \cmark & \cmark & \cmark & \cmark \\
Roth et al.~\cite{roth2022patchcore}        & \xmark & \cmark & \cmark & \cmark & \cmark \\
Zhang et al.~\cite{zhang2018lpips}          & \xmark & \cmark & \xmark & \xmark & \cmark \\
Sung et al.~\cite{sung2018relation}         & \xmark & \cmark & \cmark & \xmark & \cmark \\
\midrule
\textbf{This thesis}                        & \cmark & \cmark & \cmark & \cmark & \cmark \\
\bottomrule
\end{tabular}
\end{table}
